\chapter*{Research}
\addcontentsline{toc}{chapter}{Research}
\stepcounter{chapter}

\section{Tiredness Identifiers}
There are several different ways to identify tiredness, especially when driving. Somes examples are:
\begin{itemize}
    \item More frequent blinking – this is because your eyes and brain function to counter fatigue and discomfort
    \item Harder Blinking – prolonged opening can result in them being saw and dry, so, to combat this, your body
    makes you blink harder and more frequently to keep them lubricated
    \item Delayed reaction times – neurons in your brain slow as they process information less efficiently
    \item Your head can begin to drop as you inevitably slowly fall asleep (worst case)
    \item Yawning – this is your body’s attempt to both keep your brain awake as well as cool the body
    \item Grip loosens – your body begins to relax as you fall asleep, actions aren’t made with the same intent
    \item Mood changes – you can often become more irritable as you fall asleep
    \item Snoring – if you have already fallen asleep, you relax and this can result in snoring
\end{itemize}

\section{Driving Effected Factors}
As you get more tired, you often lose concentration. You can’t find yourself reacting more slowly and your decision
making can become impaired. You can also become a more violent or unstable driver as your mood can be easily influenced
as you do not have the energy to combat how you feel. Not only this, but in the worst case (in the event of falling
asleep) you can no longer operate the vehicle.

\section{Legality of Using a Mobile Device}
In this project, I would like to make use of the users’ own device, which will probably be their mobile phone. To get
around the legal issue of using a mobile device, it will have to be securely fastened in a cradle and not held by the
driver. Furthermore, I will not be able to request interaction from the driver once the vehicle is in motion unless I
use hands-free \cite{gov_mobile_law}, "for example:
\begin{itemize}
    \item A Bluetooth headset
    \item Voice commands
    \item A dashboard holder or mat
    \item A windscreen mount
    \item A built-in satnav” 
\end{itemize}
Failure to comply with this can result in 6 penalty points and a £200 fine. But this can go up to a driving ban as
well as a maximum of £1000 fine.

\begin{mdframed}
\textbf{Key Implementation Takeaways:}
\begin{itemize}
    \item \textbf{Setup Protocol:} User-interactive features must be restricted to the initial setup phase while the
    vehicle is stationary.
    \item \textbf{Motion Safety:} Integration of the device's internal accelerometer to disable interactive displays
    during motion, ensuring the system does not become a distraction.
\end{itemize}
\end{mdframed}

\section{Potential Models and Algorithms}
\subsection{Initial thought process – basic Convolutional Neural Netwrk (CNN):}
My initial idea was to follow a methodology I had previously learnt in one of my prior modules, using a simple
\textbf{Convolutional model} consisting of input nodes containing extracted values from an image, and then train the model to
detect these trends in data to predict an outcome. For this I made use of \textbf{Mediapipe’s} face mesh to extract features from
frames extracted from videos, to form inputs into the neural network. This was my initial thought process, and with a desire
to start coding, and deliver a prototype, I naively started to develop this without additional research.

This methodology and implementation is talked about in ‘Simple AI Models’(\ref{sec:simple-ai-models}) section of this document. Here I explain my
methodology and thought process (as well as data validation) as to why I believed this method feasible in the early stages of
this project.

Following this failure, which is also expressed in the ‘Simple AI Models’(\ref{sec:simple-ai-models}) section I expanded my research to other AI models
which will follow.

\subsection{Deep Convolutional Models:}
As an expansion from the idea above, this consists of a far greater number of layers. With its main uses in item identification within images,
this would be ideal for the task at hand. However, this adds size and complexity to the
network, which may make it more difficult to run on the Pi compared to other models due to its increased parameter.

\subsection{Xception Model (Extreme Inception):}
A development of the normal inception model, this stands as the middle ground between regular convolutional models and depthwise
separable convolutional models (depthwise + pointwise). This makes better use of parameters available compaired to the Inception v3 model.
This is done through channel mixing.

By removing the multi branch capabilities of an inception model, opting for a single branch this completely decouple spatial
correlations as well as branch correlations, which results in more focused learning .

To take this further, you can make use of pre-trained models. The example I have used in the Inception Models section of this
paper was trained on the \textbf{ImageNet} dataset. This transferred learning allows for the pre-learns identification of edges, shapes and facial
structure to be applied to my current data, without the need for the network to necessarily learn these itself, rather re-enforce its
previous learning with my new data. This should allow for a model which learns faster, requiring less time to train, as well as a more
accurate model, due to its additional pre-training.
Some information gathered from 'Xception: Deep Learning with Depthwise Separable Convolutions'\cite{xception}.

\subsection{CondenseNet:}
This is an example of an “Efficient DenseNet using Learned Group Convolutions”\cite{condensenet}. This was designed for mobile use, with the aim of
requiring limited computational power by removing connections between layers, aiming to remove redundancies which are not used in the
later layers. Due to this, only one tenth of the computation is required compared to the DenseNet (a Deep Convolutional Model where all
layers are connected to every other layer, rather than one layer being connected to the next). Comparable to MobileNet (an Inception
model approach which focuses on efficiency and power usage rather than accuracy) but requiring less computation.
This is done through pruning in the early stages of learning to form a ‘sweet spot’ \cite{condensenet} in between sparsity and
regularity. Essentially, group (cluster) sizes are reduced (e.g. $3 \times 3 \to 1 \times 1$) which results in less computation, but a
loss in accuracy and requirement of more grouping of inputs.
The idea of this is to reduce the number of connections during learning, to improve computational speeds and re-enforce routes, at the
cost of accuracy.

\subsection{Further Development Ideas / Key Takeaways:}
\paragraph{Model Selection}
On such a small number of classifications and training examples for each classification, \textbf{CondenseNet} may not be the best option for
this project. While it is said to be faster and more optimisable for a device with limited computational ability and power (like the
Raspberry Pi I am using), it requires a greater quantity of training data than I have, without crossing between sources, and for the
sake of this project, I believe that this may be beyond the scope of what is required for success – but this may change as the project
progresses.
Xception, while slower, can provide greater accuracy given the training data that I have. This would be good for experimenting on what
the greatest achievable accuracy of a model given my training data could be. Not only this, but this model is available in Tflite, so
it is possible for the performance to be improved while on the Pi. While not necessary, this is a nice bonus and something worth
considering.
Documentation on this can be found at \url{https://keras.io/api/applications/xception/}.

\paragraph{Pre-train Options:}
To improve the versatility of these AI models, the images that are used for training should be altered to allow for a more versatile
classification. 
Since users will be using this at day and night, it would be beneficial to alter some of the training data (if it isn’t already at
day and night) by adjusting the \textbf{brightness} to simulate darkness or lower light conditions such as changes in cloud cover or
going through a tunnel. \textit{(±15°)}

The camera angle and position can also change depending on how the user sets up the system. While the location of the recordings
cannot be changed in post, the \textbf{angle} can. This will help to provide some fixability in the model, to still allow for detection
of tiredness, even if the set-up is not perpendicular to the driver, like in most of the recordings. \textit{(±20\%)}

Another issue is \textbf{zoom}. My hardware, as well as the hardware between the different training data sources, is different. As a
result, my inputs may automatically be different based on camera zoom, and how the user positions their face in the focus region,
so fluctuations in zoom/crop may be a thing to consider.  (0.9 $\to$ 1.1 zoom/crop)

\subsection{Information Communication:}
\paragraph{Webpage Hosting}
With this project being done in python, for ease of use and versatility of existing libraries for the AI section of this project, it
would be best to utilise some existing python libraries for the webpage.

\subparagraph{Flask and Alternate Frameworks}
Flask is a lightweight web framework used for web-based interfaces and API’s. In the case of this project, it would be used to host
the local webpage. This will allow for offline communication of the drowsiness state. 

\textbf{Tkinter} and \textbf{PyQt} provide a graphical user interface (idea for what I need), however, have a far more limited flexibility as well as
requiring a physically connected device. This will further limit how the project will function and be operated, as additional wiring
will create clutter. As a result, neither of these frameworks are likely to be considered in the final result.

\textbf{Django} is a very powerful framework, however, with this comes complexity for a system with such simple routing and visualisation
requirements.

Another framework to consider is \textbf{REACT}. With past experience in other projects, this has the upside that I understand
it already. However, with the project being aimed to be as light as possible to improve efficiency, this client-heavy framework might
not be the best for this project.

Overall, the low resource requirements and lightweight nature of this project, as well as the majority of the code base being in
python, the obvious choice seems to be \textbf{Flask} to provide a platform-independent web interface.

\subparagraph{API}
A simple way of connecting the webpage to the AI model. There are several different ways of communicating data from the tiredness
detection AI to the webpage:
\begin{itemize}
    \item \textbf{Shared File:} The idea is that both programs for AI detection and web hosting are run simultaneously with a common
    file to relay information between them. However, this could be more compute heavy, something I have expressed repeatedly as being
    something I want to reduce. Additionally, this can result in race conditions, where there is inconsistency between the file and the
    actual value. From a hardware point of view, with files being stored on the SD card (with limited numbers of writes) this method
    can result in shortfalls in the long run.
    \item \textbf{Shared Database:} Using some form of database (e.g. SQLite) can provide additional information as well as filtering.
    However, because of the simplistic nature of this project, it is likely too costly to do this. While given more time, and a greater
    fixation on tracking and suggested improvements to inform the user of patterns or issues, this is not in the scope of this project. 
    \item \textbf{Use of integrated API:} This is the least expensive option in terms of computing power. You can make use of threading
    to keep the two lightweight, related programs, together. This makes it easier to manage and in the event of failure of the AI, the
    whole program knows (as opposed to being in the dark due to each program being separate).
\end{itemize}

Therefore, \textbf{Option 3} (combined server and AI runner with threading) was selected as it minimises computational overhead, avoids
disk I/O bottlenecks, and keeps the architecture simple for this real-time detection application.